% Active characters and single-character control sequences are separate
% namespaces, not the same one.
%
% tex.web keeps them in genuinely separate eqtb regions (active_base vs
% single_base): the active character `~` and the control sequence `\~` are
% two different things, and \catcode is what decides which one a bare
% character refers to. Collapsing both into the same string key -- as this
% engine used to, via the bare character as its own name -- let plain.tex's
% own bootstrap corrupt a real primitive: line 526's
% `{\obeyspaces\global\let =\space}` defines the *active* space
% \obeyspaces makes, and silently overwrote the pdfTeX primitive `\ `
% (control-space, tex.web's ex_space) purely because both were keyed on the
% string " ". Every `\ ` in every document after that point meant
% `\space` instead of the real primitive.
%
% `\ ` was missing outright once the collision was fixed and stopped
% masking it -- it is added here as its own primitive, verified to behave
% exactly like an ordinary space character (both in paragraph-starting and
% in the glue it produces).
%
% Also verified: \let copying an active character's *current meaning* (not
% treating it as a literal character to bind against), and \ifx comparing
% an active character against a control sequence by looked-up meaning.
\input plain
\def\M#1{\message{<#1>}}
\catcode`\~=13
\def~{TILDE}
\M{a=~}
\let\foo~
\M{b=\foo}
\M{c=\ifx~\foo Y\else N\fi}
\let\bar=\relax
\M{d=\ifx~\bar Y\else N\fi}
\def\samething{TILDE}
\M{e=\ifx~\samething Y\else N\fi}
\catcode`\_=11
\M{f=\ifx~\undefinedcs Y\else N\fi}
\bye
